\chapter{اللقاء التاسع والعشرون: نكاح المشركات وآيات المحيض والحرث}

\section{مقدمة}
السلام عليكم ورحمة الله وبركاته. الحمد لله رب العالمين، وأشهد أن لا إله إلا الله وحده لا شريك له، وأشهد أن محمداً عبده ورسوله، صلى الله عليه وعلى آله وصحبه وسلم.

استأنف الشيخ هذا اللقاء من قوله تعالى: \begin{ayah}وَلَا تَنْكِحُوا الْمُشْرِكَاتِ حَتَّى يُؤْمِنَّ\end{ayah}، ثم مضى في آيات المحيض وإتيان النساء وآداب ذلك، حتى ختم بالكلام على قوله تعالى: \begin{ayah}وَاتَّقُوا اللَّهَ وَاعْلَمُوا أَنَّكُمْ مُلَاقُوهُ وَبَشِّرِ الْمُؤْمِنِينَ\end{ayah}. وكان هذا الموضع من أنفس المواضع في دراسة العموم والخصوص، والتخصيص والنسخ، ودلالة الحروف، وربط التفسير بنظائر القرآن.

\separator

\section{تأويل (ولا تنكحوا المشركات حتى يؤمن)}
\isnad{قال أبو جعفر:} اختلف أهل التأويل في هذه الآية:
\begin{itemize}
    \item فقيل: المراد كل مشركة، ثم نسخ نكاح الكتابيات بآية المائدة.
    \item وقيل: إن الآية من أول الأمر خاصة بمشركات العرب ومن ليس من أهل الكتاب، فلا نسخ فيها.
    \item وقيل: بل هي عامة في كل مشركة، ولا يدخل فيها استثناء.
\end{itemize}

ثم رجح \rawi{الطبري} قول \rawi{قتادة}: أن المراد بالمشركات هنا من لم يكن من أهل الكتاب، وأن الآية عام ظاهرها خاص باطنها، فلا نسخ فيها، وإنما هي من باب العموم الذي أريد به بعض أفراده.

ونبّه الشيخ هنا إلى دقة التفريق بين ثلاثة أمور متقاربة قد تختلط على القارئ:
\begin{itemize}
    \item العموم الباقي على عمومه.
    \item والعموم الذي دخله تخصيص.
    \item والعموم الذي أريد به الخصوص من أول الأمر.
\end{itemize}

وزاد أن هذا الموضع من أنفع ما يدرب طالب العلم على حسن تصوير محل النزاع قبل الترجيح؛ لأن الخلاف هنا ليس في الحكم فقط، بل في نوع العموم نفسه: أيراد به جميع أفراده ثم استثني بعضهم، أم أن بعض الأفراد لم يدخل فيه ابتداء؟

\begin{fawaid}[ليس كل تعارض ظاهري بين الآيات يوجب القول بالنسخ]
من القواعد النفيسة عند الطبري أن الحكم على آية بأنها ناسخة لأخرى لا يجوز بمجرد الدعوى، بل لا بد من برهان قاطع. فإذا أمكن حمل النصين على وجه صحيح من التخصيص أو البيان أو العموم المراد به الخصوص، كان ذلك أولى من المصير إلى النسخ.
\end{fawaid}

\section{قاعدة الطبري في النسخ}
قرر \rawi{الطبري} هنا قاعدة مهمة جداً: أن كل آيتين أو خبرين بدا في الظاهر أن أحدهما نفي حكم الآخر، فلا يجوز الجزم بأن أحدهما ناسخ للآخر إلا بدليل قاطع يزيل الاحتمال. أما مجرد الدعوى فلا وزن لها، ولهذا قال بمعنى جليل: إن المدعي بلا برهان متحكم، والتحكم لا يعجز عنه أحد.

ومن تمام صنيعه هنا أنه لم يكتف برد دعوى النسخ، بل بيّن أن آية المائدة يمكن أن تكون مبينة لمراد آية البقرة وكاشفة لخصوصها، لا رافعة لحكم ثبت ثم نسخ. وهذا هو المسلك الذي يليق بما عرف من تضييقه باب النسخ.

\begin{fawaid}[النسخ لا يثبت بالدعوى المجردة]
من أعظم ما يحفظ باب التفسير والفقه من الفوضى أن لا يدعى النسخ إلا ببينة. فكل من رأى نصين وبينهما ظاهر تعارض لا يجوز له أن يبادر إلى النسخ ما دام الجمع أو التخصيص أو البيان ممكناً.
\end{fawaid}

\separator

\section{تأويل (والمحصنات من المؤمنات والمحصنات من الذين أوتوا الكتاب)}
من خلال هذا الموضع قرر \rawi{الطبري} أن آية المائدة ليست ناسخة لآية البقرة عند اختياره، بل هي قرينة على أن الكتابيات لم يدخلن ابتداء في عموم \textit{المشركات} المراد تحريمهن هنا.

وفي ذلك تدريب لطالب العلم على أن بعض الآيات يفسر بعضاً، وأنه قد يستفاد من آية متأخرة تقرير معنى آية سابقة من غير لزوم النسخ.

كما ظهر هنا معنى آخر مهم في صناعة الترجيح: أن الأثر الذي يخالف ظاهر القرآن المستقر أو إجماع الأمة لا يقوى على معارضته، ولذلك لم يجعل \rawi{الطبري} ما روي في المنع عن بعض الصحابة معارضاً للحكم الثابت، بل حمل الثابت عن \rawi{عمر} رضي الله عنه على معنى الكراهة السياسية أو الاحتياطية لا على التحريم.

\separator

\section{تأويل (ولا تنكحوا المشركين حتى يؤمنوا)}
كما حرم الله على المؤمنين نكاح المشركات حتى يؤمن، حرم على المؤمنات أن يبقين أو يزوجن للمشركين حتى يؤمنوا. وفي هذا تقرير لمعيار عظيم: أن رابطة الإيمان مقدمة على اعتبارات النسب والجاه والهوى.

\section{تأويل (ولأمة مؤمنة خير من مشركة ولو أعجبتكم)}
بيّن \rawi{الطبري} أن ميزان الاختيار في النكاح ليس الجمال المجرد، ولا الحسب، ولا المال، بل الإيمان الذي يورث البركة والصلاح. وهذا المعنى يعم الرجال والنساء جميعاً.

\separator

\section{تأويل (ويسألونك عن المحيض)}
\isnad{قال أبو جعفر:} المحيض هو الدم المعروف أو زمانه وموضعه بحسب السياق، والمقصود هنا النهي عن قربان النساء في حال الحيض في موضع الحرث.

ثم أوضح \rawi{الطبري} أن الله لم ينههم عن مؤاكلة الحائض ولا مساكنتها ولا سائر وجوه المعاشرة المباحة، وإنما نهى عن الجماع في زمن الحيض. وهذا من تصحيح الإسلام لغلو الجاهلية وأهل الكتاب في باب الحائض.

\begin{fawaid}[الشريعة ترفع الغلو كما ترفع التفريط]
من محاسن القرآن أنه لا يكتفي بالنهي، بل يصحح المفهوم المحيط بالحكم. ففي باب الحيض لم يقر غلو من كان يهجر الحائض هجراً تاماً، ولم يفتح باب ما حرم الله، بل توسط بالعدل والبيان.
\end{fawaid}

\section{تأويل (فإذا تطهرن فأتوهن من حيث أمركم الله)}
ذكر \rawi{الطبري} الأقوال في معنى \textit{من حيث أمركم الله}، ثم رجح أن المراد: فأتوهن في الفروج من الوجه الذي أذن الله لكم بإتيانهن منه، وهو القبل، بعد الطهر والتطهر، لا في حال الحيض ولا في غير موضع الحرث.

وبين الشيخ أن هذا الموضع من أجمل أمثلة رد التفسير الفاسد بجمع القرآن إلى السنة إلى اللسان العربي. فليس كل معنى محتمل في اللفظ يصح حمل الآية عليه، ما دام يخالف نصوصاً أوضح منه.

\begin{fawaid}[القرآن يفسر بالسنة وباللسان معاً]
منهج الطبري هنا بديع: لا يكتفي باللسان وحده، ولا بالأثر مجرداً، بل يجمع بين دلالة القرآن، والسنة الثابتة، ولسان العرب، حتى يستقيم المعنى وتزول الشبهات.
\end{fawaid}

\section{تأويل (إن الله يحب التوابين ويحب المتطهرين)}
ذكر \rawi{الطبري} أقوالاً في معنى \textit{المتطهرين}:
\begin{itemize}
    \item فقيل: المتطهرون بالماء للصلاة.
    \item وقيل: المتطهرون من إتيان النساء في الأدبار.
    \item وقيل: المتطهرون من الذنوب فلا يعودون إليها.
\end{itemize}

ثم رجح أن المراد أعم ما يظهر من السياق: التطهر بالماء من الأحداث والحيض والجنابة ونحوها، مع ما في اللفظ من إشعار بمحبة الله لعباده إذا رجعوا إليه وتطهروا ظاهراً وباطناً.

\separator

\section{تأويل (نساؤكم حرث لكم فاتوا حرثكم أنى شئتم)}
\isnad{قال أبو جعفر:} النساء موضع الحرث، أي: موضع الزرع والنسل، ولذلك سمين حرثاً من جهة أن الأولاد يخرجون منهن كما يخرج الزرع من موضعه.

ثم تكلم \rawi{الطبري} على معنى \textit{أنى}، وبيّن أنها ليست بمعنى \textit{متى} ولا \textit{أين} ولا \textit{كيف} على الإطلاق، بل هي أدق من ذلك، وتدل هنا على جهات الإتيان ووجوهه في موضع الحرث المشروع. وميزة بحثه هنا أنه لم يكتف بالدعوى اللغوية، بل فرّق بين هذه الأدوات من جهة الأجوبة التي تقتضيها العرب في كلامها، فجعل ذلك من أقوى ما يرد به التفسير الفاسد.

ومن ثم رجح أن المعنى: ائتوا نساءكم في القبل من أي جهة شئتم: مقبلات أو مدبرات أو باركات، لا أن الآية إباحة للإتيان في الدبر؛ لأن الدبر ليس موضع حرث أصلاً.

\begin{fawaid}[تحرير دلالة الحرف الواحد قد يغير فهم الآية كلها]
من أعظم ما يظهر به عمق التفسير معرفة معاني الأدوات والحروف. فبحث الطبري في معنى \textit{أنى} هنا ليس ترفاً لغوياً، بل هو مفتاح الحكم كله؛ إذ به يتبين أن الآية في جهات الإتيان لا في تغيير موضعه المشروع.
\end{fawaid}

\section{سبب النزول وأثره في الترجيح}
استوعب \rawi{الطبري} ما روي في سبب نزول الآية من اعتراض اليهود على بعض صور إتيان الرجل امرأته في القبل إذا كان من جهة دبرها، ثم جعل هذا من أقوى ما يفسر به الآية، لأن السبب يكشف السؤال الذي نزل الجواب لأجله.

\begin{fawaid}[سبب النزول يضبط جهة الجواب]
إذا عرفت السؤال الذي نزلت الآية جواباً له، انكشف لك كثير من الاحتمالات البعيدة. ولذلك كان سبب النزول هنا قرينة قوية على أن الآية في هيئة الإتيان في القبل، لا في إباحة موضع محرم.
\end{fawaid}

\section{تأويل (وقدموا لأنفسكم)}
ذكر \rawi{الطبري} في هذه الجملة قولين مشهورين:
\begin{itemize}
    \item فقيل: قدموا ذكر الله عند الجماع.
    \item وقيل: قدموا لأنفسكم الخير والطاعة والعمل الصالح.
\end{itemize}

ثم رجح الثاني، واستدل له بآية أخرى: \begin{ayah}وَمَا تُقَدِّمُوا لِأَنْفُسِكُمْ مِنْ خَيْرٍ تَجِدُوهُ عِنْدَ اللَّهِ\end{ayah}. ففسر القرآن بالقرآن، وجعل الأمر هنا عاماً في كل ما يقدمه العبد لنفسه من خير، وإن دخل فيه ذكر الله عند الجماع دخولاً أولياً أو تبعياً.

ونبّه الشيخ إلى فائدة دقيقة في هذا الترجيح: أن الطبري لم ينكر دخول بعض الأفراد المذكورة في الآية، لكنه أبى حصر اللفظ العام فيها مع وجود شاهد قرآني أوضح يرده إلى معنى أوسع وأجمع.

\begin{fawaid}[من أجود وجوه التفسير تفسير القرآن بالقرآن]
إذا وجد في القرآن لفظ قد فسر استعماله في موضع آخر أوضح، كان حمله على ذلك من أقوى وجوه الترجيح. وقد أحسن الطبري هنا حين جعل آية البقرة تفسرها آية أخرى من القرآن نفسه.
\end{fawaid}

\separator

\section{تأويل (واتقوا الله واعلموا أنكم ملاقوه وبشر المؤمنين)}
ختم \rawi{الطبري} هذا المقطع بأصول جامعة:
\begin{itemize}
    \item التقوى: في امتثال الأحكام السابقة واجتناب محارم الله فيها.
    \item والعلم باللقاء: لأنه المحرك الأعظم للتقوى والاستقامة.
    \item والبشارة: لأن من أحسن العمل وثبت على الإيمان فحقه أن يبشر برحمة الله ورضوانه.
\end{itemize}

ونبّه الشيخ إلى أن البشارة للعاملين من أعظم ما يثبتهم ويحبب إليهم الطاعة، وأن من التربية أن يثنى على الخير ويذكر صاحبه بثوابه، لا أن لا يلتفت إليه إلا عند الخطأ.

\begin{fawaid}[البشارة بالخير من التربية الإيمانية]
كما أن التخويف يزجر النفوس، فإن البشارة تثبتها وتحبب إليها الطريق. ولذلك كان من تمام التربية أن يبشر العامل بثواب عمله، وأن يذكر المؤمن بما أعد الله له من الكرامة.
\end{fawaid}

\separator

ختم الشيخ اللقاء بشكر الله على تمام المجلد الثالث، والتأكيد على أن هذه المقاطع تظهر جمال منهج \rawi{الطبري} في دراسة النصوص النقدية: حسن تصوير المسألة، وجمع الأقوال، وتمييز العموم من الخصوص، والنسخ من التخصيص، والاحتمال الصحيح من الفاسد، وربط ذلك كله بالقرآن والسنة ولسان العرب.
